\chapter{Experimental Results}
\label{chap:experimental-results}

\section{Measurement methodology and scopes}

This chapter reports the gas and runtime measurements obtained during the experimental
phase of the project. The main difficulty was not only measuring gas, but measuring the
same quantity across different architectures. The implementation went through three
relevant stages:

\begin{itemize}
  \item the initial implementation, used as the baseline from the previous semester
  project;
  \item the intermediate monolithic implementation, where several variants were first added to
  the existing contracts as flags;
  \item the final specialized implementation, where the variants are split into dedicated
  contracts and, for the optimistic variants, into reusable implementations deployed through
  a factory.
\end{itemize}

All gas values in this chapter come from transaction receipts in the local Hardhat test
environment. They should therefore be read as reproducible EVM gas measurements for the
implemented contracts, not as predictions of end-user latency or fee paid on a public
network. Off-chain timings were measured with the native Rust command-line tools used by
the benchmarks. They are machine-dependent, but they are useful for comparing the relative
cost of preparing precontracts and dispute witnesses.

The measurements were obtained from the Phase 4 benchmark suite, in particular the tests
comparing the initial implementation, the monolithic implementation, the specialized final
architecture, and the native 1 GiB hardcoded dispute path. The benchmark outputs were then
reorganized by scope before being reported here, so that totals can be recomposed from the
individual lines when needed.

The following scopes are used throughout the chapter.

\begin{table}[htbp]
  \centering
  \begin{tabularx}{\linewidth}{>{\raggedright\arraybackslash}p{0.27\linewidth}>{\raggedright\arraybackslash}X}
    \toprule
    Scope & Meaning \\
    \midrule
    Fixed infrastructure &
    Gas paid once to deploy reusable components such as \texttt{SOXFactory} and optimistic
    implementation contracts. \\
    Marginal optimistic cost &
    Gas paid for one new exchange after the reusable infrastructure already exists. For
    clone-based modes, this includes clone creation and the honest-path calls until
    completion. \\
    Optimistic path &
    Gas needed to deploy or create the optimistic account and execute the honest exchange
    until the terminal complete state. \\
    Dispute after trigger &
    Gas paid after the dispute account has been deployed, namely the challenge rounds,
    Step 8 submission, and finalization. \\
    Full dispute path &
    Gas including the optimistic setup, dispute trigger/deployment, challenge rounds,
    Step 8, and finalization. \\
    Off-chain runtime &
    Wall-clock time for native computation of precontract artifacts or dispute witnesses.
    It is not on-chain gas. \\
    \bottomrule
  \end{tabularx}
  \caption{Measurement scopes used in the experimental evaluation.}
  \label{tab:measurement-scopes}
\end{table}

This separation is essential. For example, a factory-based clone has a very low marginal
cost per exchange, but only after the factory and implementations have already been
deployed. Conversely, a direct hardcoded deployment can be more expensive in the optimistic
phase while still being useful in dispute, because it reduces the dispute deployer and the
generic circuit verification work.

\section[Baseline measurements]{Baseline measurements of the initial implementation}

The first baseline is the honest optimistic exchange in the initial implementation. In that
architecture, every exchange deploys a full optimistic account contract. The complete
optimistic path measured in the initial implementation costs \(2,300,201\) gas: \(2,077,362\)
gas for deployment and \(222,839\) gas for execution.

\begin{table}[htbp]
  \centering
  \begin{tabular}{lrrr}
    \toprule
    Optimistic path & Initial & Interm. mono. & Difference \\
    \midrule
    Deployment & 2,077,362 & 2,813,862 & +736,500 \\
    Execution until complete & 222,839 & 232,920 & +10,081 \\
    \midrule
    Total & 2,300,201 & 3,046,782 & +746,581 \\
    \bottomrule
  \end{tabular}
  \caption{Honest optimistic path before and after the first monolithic variant changes.}
  \label{tab:baseline-optimistic}
\end{table}

Table~\ref{tab:baseline-optimistic} shows the negative result that motivated the final
architecture. The monolithic implementation was functionally useful, but it made the
optimistic path \(32.46\%\) more expensive than the initial implementation. Most of the
increase came from deployment, not from the honest-path calls. This confirmed that putting
all variants into the same deployed bytecode was the wrong optimization strategy.

The same effect appears in the dispute initialization measurements on a 4 MiB file. The
monolithic implementation is again more expensive than the initial implementation when the
same scope is measured.

\begin{table}[htbp]
  \centering
  \small
  \begin{tabularx}{\linewidth}{>{\raggedright\arraybackslash}Xrrr}
    \toprule
    Dispute-related measure, 4 MiB & Initial & Interm. mono. & Difference \\
    \midrule
    Dispute deployment / trigger & 4,651,201 & 5,779,296 & +1,128,095 \\
    Exchange + dispute triggering & 4,977,003 & 6,058,222 & +1,081,219 \\
    Optimistic deploy + dispute init total & 7,054,365 & 8,872,096 & +1,817,731 \\
    One challenge round average & 82,624 & 103,282 & +20,658 \\
    \texttt{submitCommitment}, SHA-256 gate & 320,000 & 379,686 & +59,686 \\
    \bottomrule
  \end{tabularx}
  \caption{Initial implementation versus intermediate monolithic implementation on the same 4 MiB dispute scope.}
  \label{tab:baseline-dispute-4mib}
\end{table}

The baseline conclusion is therefore clear: the first implementation of variants did not
optimize gas. It was an intermediate engineering step that made the variants executable and
measurable, but its gas cost justified the move to specialized contracts.

\section{One-time infrastructure deployment costs}

The final optimistic architecture uses a reusable factory and reusable implementation
contracts. These costs are paid once at protocol setup time. They are not paid again for
each exchange.

\begin{table}[htbp]
  \centering
  \begin{tabularx}{\linewidth}{>{\raggedright\arraybackslash}Xr}
    \toprule
    Component & Deployment gas \\
    \midrule
    \texttt{SOXFactory} & 561,126 \\
    Normal optimistic clone implementation & 1,492,453 \\
    \texttt{no\_S\_deposit} implementation & 1,482,093 \\
    \(S=B\) implementation & 1,529,267 \\
    \(S=V\) implementation & 1,519,062 \\
    \midrule
    Total optimistic clone infrastructure & 6,584,001 \\
    \bottomrule
  \end{tabularx}
  \caption{One-time infrastructure cost of the clone-based optimistic architecture.}
  \label{tab:factory-infrastructure-cost}
\end{table}

This table should not be added to every exchange. It is an amortized infrastructure cost.
For a single isolated exchange, deploying this infrastructure would be expensive. For a
system handling several exchanges, however, the marginal cost becomes the relevant measure.
Compared with the initial optimistic total of \(2,300,201\) gas per exchange, the clone
infrastructure is amortized after roughly four exchanges in the measured specialized
optimistic modes.

\section{Optimistic path gas costs}

The final implementation provides two kinds of specialized optimistic measurements. The
first kind uses direct specialized contracts. These already reduce the cost compared with
the initial implementation, because each contract removes unrelated branches. The second
kind uses minimal clones through \texttt{SOXFactory}. These are the lowest marginal-cost
measurements, because the full implementation bytecode has already been deployed once.

\begin{table}[htbp]
  \centering
  \small
  \begin{tabularx}{\linewidth}{>{\raggedright\arraybackslash}Xrrr}
    \toprule
    Specialized optimistic mode & Direct total & Clone marginal total & Clone gain vs initial \\
    \midrule
    \texttt{no\_S\_deposit} & 1,724,251 & 500,720 & \(-78.23\%\) \\
    \(S=B\) & 1,751,300 & 516,349 & \(-77.55\%\) \\
    \(S=V\) & 1,787,289 & 556,309 & \(-75.81\%\) \\
    \bottomrule
  \end{tabularx}
  \caption{Optimistic-path gas costs for the final specialized modes.}
  \label{tab:optimistic-specialized-costs}
\end{table}

The clone-based \texttt{no\_S\_deposit} mode is the cheapest measured honest exchange:
\(500,720\) gas once the infrastructure exists. The \(S=B\) and \(S=V\) cases are slightly
more expensive, but they remain below \(560,000\) gas. These values are the main positive
result on the optimistic path. The project reduced the per-exchange optimistic cost from
roughly \(2.3\) million gas in the initial implementation to roughly \(0.5\)--\(0.56\)
million gas in the final specialized clone architecture.

It is important not to compare these clone marginal values directly with a first deployment
of the whole protocol. The fair comparison is either direct-specialized versus initial
single-exchange deployment, or clone marginal versus initial per-exchange cost after the
clone infrastructure has been paid once. Appendix~\ref{app:detailed-gas-tables} gives an
explicit recomposition table for these final factory/clone costs.

\section{Dispute path gas costs}

For the dispute path, the main same-scope comparison uses a 4 MiB file and measures
``optimistic deploy + dispute init total''. This scope combines the cost of creating the
optimistic exchange with the cost of reaching the dispute initialization point. It is useful
because it exposes both the deployment effect and the dispute trigger effect.

\begin{table}[htbp]
  \centering
  \begin{tabularx}{\linewidth}{>{\raggedright\arraybackslash}Xr}
    \toprule
    Architecture, same 4 MiB scope & Gas \\
    \midrule
    Initial implementation & 7,054,365 \\
    Intermediate monolithic implementation & 8,872,096 \\
    Final specialized implementation & 5,495,184 \\
    \bottomrule
  \end{tabularx}
  \caption{Same-scope 4 MiB dispute initialization comparison.}
  \label{tab:same-scope-dispute-4mib}
\end{table}

The final specialized implementation reduces this scope by \(38.06\%\) compared with the
monolithic implementation and by \(22.10\%\) compared with the initial implementation. The
main reason is architectural: the optimistic deployment component becomes a clone creation,
while the generic Step 8 verifier remains close to the monolithic cost in the normal
non-hardcoded case.

A second same-scope comparison uses 16 KiB disputes and compares the intermediate monolithic
diagnostic against the final specialized contracts for the four relevant dispute variants.
This scope starts at the buyer-side dispute sponsor step and includes the trigger,
challenge rounds, Step 8, and finalization.

\begin{table}[htbp]
  \centering
  \small
  \begin{tabularx}{\linewidth}{>{\raggedright\arraybackslash}Xrrr}
    \toprule
    Dispute case, 16 KiB & Interm. mono. & Final specialized & Gain \\
    \midrule
    Normal, external sponsors & 8,345,424 & 7,432,647 & \(-10.94\%\) \\
    Self-sponsored \(S_B=B, S_V=V\) & 7,165,928 & 5,946,445 & \(-17.02\%\) \\
    Hardcoded \(\mathsf{SHA256}\) & 8,310,230 & 5,195,022 & \(-37.49\%\) \\
    Self-sponsored + hardcoded & 7,193,323 & 4,042,185 & \(-43.81\%\) \\
    \bottomrule
  \end{tabularx}
  \caption{Same-scope 16 KiB dispute comparison between intermediate monolithic and final specialized contracts.}
  \label{tab:same-scope-dispute-16kib}
\end{table}

Table~\ref{tab:same-scope-dispute-16kib} is the clearest dispute-side summary of the final
architecture. The normal case improves moderately because the verifier is still generic.
The hardcoded cases improve much more because the dispute deployer and Step 8 verifier are
specialized. The self-sponsored hardcoded case is the best measured 16 KiB dispute case in
this scope, with a total of \(4,042,185\) gas.

\section{Hardcoded SHA-256 gas impact}

The hardcoded \(\mathsf{SHA256}\) mode has a specific purpose. It does not optimize the
honest optimistic path directly. It optimizes the dispute path for the common description
\(desc=\mathsf{SHA256}(x)\), by making the smart contract reconstruct the selected gate
from public metadata instead of receiving the generic circuit gate and the Merkle proof
\(\pi_1\) for \(h_{\mathit{circuit}}\).

The local Step 8 measurement confirms that \(\pi_1\) is removed. For a 16 KiB circuit, the
ordinary \texttt{submitCommitment} call costs \(463,532\) gas and carries 10 \(\pi_1\)
items. The hardcoded call costs \(400,990\) gas and carries no \(\pi_1\) item, saving
\(62,542\) gas on that local Step 8 call.

\begin{table}[htbp]
  \centering
  \begin{tabular}{lrrrr}
    \toprule
    Size & Normal Step 8 & Hardcoded Step 8 & \(\pi_1\) items & Saving \\
    \midrule
    13 bytes & 251,444 & 234,626 & 3 to 0 & 16,818 \\
    16 KiB & 463,532 & 400,990 & 10 to 0 & 62,542 \\
    \bottomrule
  \end{tabular}
  \caption{Local Step 8 effect of removing the generic \(h_{\mathit{circuit}}\) proof.}
  \label{tab:hardcoded-step8-local}
\end{table}

For large circuits, isolating only the \(h_{\mathit{circuit}}/\pi_1\) membership check gives
a larger local saving. At 1 GiB, the generic membership verification costs approximately
\(198,585\) gas, while the hardcoded reconstruction costs between \(29,639\) and \(32,390\)
gas depending on the selected gate. This is a local saving of about \(166,000\)--\(169,000\)
gas. In a full dispute, this saving is diluted by the other costs: dispute deployment,
challenge rounds, AES/SHA gate evaluation, \(h_{\mathit{ct}}\) proofs, \(hpre\) proofs, and
finalization.

\begin{table}[htbp]
  \centering
  \small
  \begin{tabularx}{\linewidth}{>{\raggedright\arraybackslash}Xrrr}
    \toprule
    Equivalent circuit size & Generic \(h_{\mathit{circuit}}\) check & Hardcoded check & Local saving \\
    \midrule
    900 MiB & 192,267 & 29,639--32,390 & 159,877--162,628 \\
    1 GiB & 198,585 & 29,639--32,390 & 166,195--168,946 \\
    \bottomrule
  \end{tabularx}
  \caption{Isolated large-circuit effect of replacing \(\pi_1\) by hardcoded reconstruction.}
  \label{tab:hardcoded-large-local}
\end{table}

The largest full-path gain of the hardcoded work comes from the specialized hardcoded
dispute deployer. In the intermediate monolithic hardcoded path, triggering a hardcoded
dispute still paid for large generic deployment logic. In the final specialized hardcoded
path, the trigger/deploy component decreases from \(5,767,448\) gas to \(2,699,433\) gas in
the measured deployment-and-trigger scope. This explains why the final hardcoded dispute
numbers improve far more than the isolated \(\pi_1\) saving alone. A full component-level
mirror of the monolithic hardcoded path is given in Appendix~\ref{app:detailed-gas-tables}.

\section{Self-sponsoring gas impact}

Self-sponsoring affects the dispute path differently from hardcoded \(\mathsf{SHA256}\).
It does not change the circuit verifier. Instead, it changes the number of times the
challenge loop must be executed. In the normal external-sponsor case, Step 9 may require a
new loop with another sponsor configuration. In the \(S_B=B, S_V=V\) case, this relaunch is
unnecessary, because the disputing parties are already their own dispute sponsors.

The 16 KiB monolithic diagnostic makes this effect visible if one removes the fixed trigger
and setup components and looks at the variable dispute part after the trigger. The normal
external-sponsor case costs \(2,450,528\) gas after the trigger. The self-sponsored case
costs \(1,277,852\) gas after the trigger, a reduction of \(47.85\%\). The self-sponsored
hardcoded case costs \(1,215,298\) gas after the trigger, a reduction of \(50.41\%\) against
the normal variable part.

This is consistent with the protocol expectation: the gas reduction is not a tiny local
saving, but comes from removing a repeated Step 7--8 loop. In full totals, the percentage
is smaller because fixed costs such as dispute deployment remain present. This distinction
between full totals and variable loop cost was one of the main lessons of the measurement
phase.

\section{Large-file measurements up to 1 GiB}

The implementation was also tested on a 1 GiB file using the native Rust path and the final
hardcoded \(\mathsf{SHA256}\) dispute. This is the most important scalability measurement of
the project: it demonstrates that the implementation can generate the required off-chain
artifacts and execute the complete on-chain dispute path up to the terminal state.

\begin{table}[htbp]
  \centering
  \small
  \begin{tabularx}{\linewidth}{>{\raggedright\arraybackslash}Xrr}
    \toprule
    Final hardcoded strict path & 16 KiB & 1 GiB \\
    \midrule
    Optimistic before trigger & 2,938,363 & 2,938,363 \\
    Trigger dispute / deploy dispute & 2,721,171 & 2,721,171 \\
    Total optimistic + trigger & 5,659,534 & 5,659,534 \\
    Challenge responses total & 608,190 & 1,581,282 \\
    Opinions total & 421,615 & 1,095,903 \\
    Step 8 \texttt{submitCommitmentLeft} & 5,449,017 & 5,577,706 \\
    Finalization & 75,667 & 75,667 \\
    Dispute after trigger & 6,554,489 & 8,330,558 \\
    \midrule
    Complete measured path & 12,214,023 & 13,990,092 \\
    \bottomrule
  \end{tabularx}
  \caption{Complete final hardcoded strict path for 16 KiB and 1 GiB.}
  \label{tab:hardcoded-strict-end-to-end}
\end{table}

The 1 GiB file has \(16,777,216\) ciphertext blocks and \(33,554,437\) gates in the
hardcoded layout. The binary search therefore uses 26 challenge rounds. In the measured
self-sponsored hardcoded scenario, the complete path costs \(13,990,092\) gas, including
the optimistic path, the dispute trigger, all challenge rounds, Step 8, and finalization.

A separate 1 GiB hardcoded benchmark with external sponsors and three Step 7--8 iterations
was also run. The dispute after trigger costs \(23,865,946\) gas. Including the trigger
gives \(26,585,279\) gas, and including the full optimistic path before the dispute gives
\(29,532,277\) gas. This is not a simple multiplication by three of the one-iteration cost,
because later iterations rewrite some storage slots that have already been initialized.

\section{Off-chain runtime measurements}

The 1 GiB benchmark uses the native Rust CLI, not the browser/WASM path. This distinction is
important. The WASM path remains useful for integration and small or medium files, but the
native CLI is the reliable path for large-file measurements because it avoids the memory
pressure of JavaScript and browser-adjacent execution.

\begin{table}[htbp]
  \centering
  \begin{tabularx}{\linewidth}{>{\raggedright\arraybackslash}Xr>{\raggedright\arraybackslash}X}
    \toprule
    1 GiB runtime measure & Time & Interpretation \\
    \midrule
    Precontract CLI & 22.42--48.70 s & Mainly vendor-side preparation: read \(x\), encrypt,
    compute commitments and metadata. \\
    Native dispute generation & 16.43--35.46 s & Aggregate benchmark time to prepare dispute
    data; not attributable to one protocol actor. \\
    Total \(hpre_i\) generation & 9.47--11.92 s & Buyer-side generation of the published
    challenge responses over 26 rounds. \\
    Local Hardhat transaction execution & 0.30--0.56 s & Local wall-clock time to send the
    benchmark transactions; not a cryptographic runtime. \\
    \bottomrule
  \end{tabularx}
  \caption{Off-chain runtime measurements for the 1 GiB hardcoded benchmark.}
  \label{tab:offchain-runtime-1gib}
\end{table}

The ranges in Table~\ref{tab:offchain-runtime-1gib} reflect repeated runs on the same
machine. The gas values are stable transaction-receipt measurements, but wall-clock times
depend on cache state and local I/O. The first \(hpre_i\) computations are slower than the later ones because the benchmark
forces the binary search to the left. The challenge index is approximately divided by two
at each round, so the native code recomputes a smaller prefix/accumulator for later rounds.
This affects CPU time only. It is not a gas effect.

\begin{table}[htbp]
  \centering
  \begin{tabular}{rrr}
    \toprule
    Round & Challenge index \(i\) & \(hpre_i\) time \\
    \midrule
    0 & 16,777,219 & 4.58--6.04 s \\
    1 & 8,388,610 & 2.23--2.46 s \\
    2 & 4,194,305 & 1.17--1.33 s \\
    3 & 2,097,153 & 0.60--0.71 s \\
    4 & 1,048,577 & 0.31--0.39 s \\
    \bottomrule
  \end{tabular}
  \caption{First \(hpre_i\) generation times in the 1 GiB forced-left benchmark.}
  \label{tab:hpre-first-rounds}
\end{table}

The on-chain gas of the challenge rounds does not follow this division by two. For the
1 GiB benchmark, \texttt{respondChallenge} costs \(1,581,282\) gas over 26 rounds, or about
\(60,819\) gas per round. \texttt{giveOpinion} costs \(1,095,903\) gas over 26 rounds, or
about \(42,150\) gas per round.

\section{Comparison summary}

The final comparison is summarized in Table~\ref{tab:experimental-summary}. The most
important result is that the final architecture does not merely recover from the
monolithic overhead; it improves over the initial implementation on the main same-scope
measurements that were targeted by the project.

\begin{table}[htbp]
  \centering
  \small
  \begin{tabularx}{\linewidth}{>{\raggedright\arraybackslash}p{0.34\linewidth}>{\raggedright\arraybackslash}X}
    \toprule
    Result & Interpretation \\
    \midrule
    Optimistic clone modes: 500,720--556,309 gas &
    Per-exchange optimistic marginal cost is reduced by about \(75.8\%\)--\(78.2\%\)
    compared with the initial \(2,300,201\) gas optimistic path. \\
    4 MiB specialized dispute-init scope: 5,495,184 gas &
    The final specialized architecture is \(38.06\%\) cheaper than the intermediate monolithic
    version and \(22.10\%\) cheaper than the initial implementation on the same scope. \\
    16 KiB specialized dispute cases &
    All four measured variants improve over the intermediate monolithic equivalents; the best
    case is self-sponsored + hardcoded at \(4,042,185\) gas. \\
    Self-sponsoring &
    The variable dispute loop cost drops by roughly half when Step 9 no longer relaunches
    unnecessary Step 7--8 iterations. \\
    Hardcoded \(\mathsf{SHA256}\) &
    The contract reconstructs \(g_i\), removes \(\pi_1\), and reduces both local Step 8 gas
    and the specialized hardcoded dispute deployment path. \\
    1 GiB hardcoded benchmark &
    A complete hardcoded dispute was executed up to \texttt{End}, with 26 challenge rounds
    and a full measured cost of \(13,990,092\) gas in the self-sponsored hardcoded scenario. \\
    \bottomrule
  \end{tabularx}
  \caption{Summary of the main experimental results.}
  \label{tab:experimental-summary}
\end{table}

\section{Interpretation of results}

The experimental results support three conclusions.

First, specialization is necessary. The monolithic implementation made the variants
available, but it increased deployment gas and made several tables hard to interpret. The
final split architecture is measurably better because mutually exclusive logic is no longer
deployed together.

Second, the optimistic path is where the factory and clone architecture is most effective.
The reduction from about \(2.3\) million gas to about \(0.5\) million gas per exchange is
the largest clean improvement of the project. It directly addresses the cost of the honest
case, which is the path expected to be used most often.

Third, dispute optimization has two different components. Self-sponsoring reduces the
number of dispute iterations, while hardcoded \(\mathsf{SHA256}\) reduces the generic
circuit-verification burden and the hardcoded dispute deployment path. These optimizations
are complementary: the strongest 16 KiB dispute result is obtained when both are enabled.

The remaining dominant costs are \texttt{triggerDispute}, deployment of dispute accounts,
and the Step 8 evaluation of expensive gates, especially AES-related checks and the
associated commitment proofs. These are the natural next targets for a future version of
the implementation.
